% Answer keys for Chapters 24--32. This file is \input directly inside the document.
% Every number is tested by an assert script outside the repo (see the report).

\section*{Chapter 24: Linear and Integer Programming}
\begin{enumerate}[label=\textbf{24.\arabic*}]
  \item Maximize $30x_{1}+20x_{2}$ subject to $2x_{1}+x_{2}\le10$, $x_{1}+2x_{2}\le8$,
        $x\ge0$. The vertices $(0;0)$, $(5;0)$, $(0;4)$, $(4;2)$ give profit
        $0;\,150;\,80;\,160$. The continuous optimum is $(4;2)$ with profit $160$; integer
        enumeration also gives $160$, so here rounding down happens to be exact and
        the gap is $0$.

  \item For $x,y\in P$ and $\lambda\in[0,1]$:
        $A(\lambda x+(1-\lambda)y)=\lambda Ax+(1-\lambda)Ay\le\lambda b+(1-\lambda)b=b$,
        using $A$ linear and $\lambda\ge0$. On integer sets that argument collapses:
        a convex combination of two integer points need not be integer, for example $0$
        and $1$ are feasible while their midpoint is not a member of $\mathbb{Z}^{n}$.

  \item Profit $x_{1}+2x_{2}$ at the vertices: $0$; $4$; $6$; $6$. The optima $(0;3)$ and
        $(3;1{,}5)$ both equal $6$, and the whole segment $x_{1}+2x_{2}=6$ is optimal because
        the objective is constant there; the midpoint $(1{,}5;\,2{,}25)$ is not a vertex.
        The dual is $y=(0;1)$: the price of an hour is $1$, of plywood $0$.

  \item Initial table with basis $s_{1},s_{2},s_{3}=(4;3;3)$. First pivot: $x_{1}$ enters
        (reduced cost $3$), ratios $4$ and $3$, so $s_{2}$ leaves; value $9$.
        Second pivot: $x_{2}$ enters, $s_{1}$ leaves (ratio $1$), optimum $(3;1)$
        with $z=11$. The dual prices from the final table: $(2;1;0)$, and
        $4(2)+3(1)+3(0)=11$.

  \item The dual minimizes $10y_{1}+8y_{2}$ subject to $2y_{1}+y_{2}\ge30$,
        $y_{1}+2y_{2}\ge20$, $y\ge0$; the solution $y=(40/3;\,10/3)$ has value
        $160$, the same as the primal. Meaning: one extra kg of flour is worth
        $13{,}333$ and one extra oven hour $3{,}333$.

  \item False. Rounding down is not always feasible once there are equality
        constraints or constraints with sign $\ge$. Example: maximize $x_{1}+x_{2}$ with
        $2x_{1}+2x_{2}=5$; the continuous answer $(2{,}5;\,0)$ rounds to
        $(2;0)$, which gives $4\ne5$ and is therefore infeasible.

  \item True for the forward direction, and only for optimal pairs. Complementary
        slackness $y_{i}(b_{i}-\mathbf{a}_{i}^{\top}\mathbf{x})=0$ gives:
        $y_{i}>0$ forces that constraint to be active. The converse does not hold; at a
        degenerate point there is an active constraint with $y_{i}=0$, so a zero price
        does not mean the constraint is slack.

  \item Densities $2{,}0$; $1{,}75$; $1{,}667$; $1{,}5$. The first item fills up
        ($5$ kg, $10$), then $3/4$ of the second, total weight $8$: the LP bound
        is $15{,}25$. Enumeration gives $15$ (items $1$ and $3$). Gap $0{,}25$,
        that is $1{,}639$ percent of the bound.

  \item 
Capacity $7$: the first item fills the sack, then $1/2$ of the second, so
        the LP bound is $13{,}5$; the best integer answer is $13$ (items $1$ and
        $4$, weight $7$). The gap is $0{,}5$, i.e. $3{,}704$ percent. The gap
        changes because capacity $7$ rules out the pair $(1,3)$, whose weight
        $5+3=8$ used to fit, so the best integer answer moves to items $1$ and $4$;
        the fractional and integer optima do not move in step.

  \item $x_{v}=1/2$ for every $v$ is feasible ($x_{u}+x_{v}=1$) and gives $n/2$;
        every cover must contain $n-1$ vertices, since leaving out two vertices
        means the edge between them is uncovered. Ratio of LP to integer:
        $0{,}6667$; $0{,}625$; $0{,}5556$ for $n=4,5,10$. The ratio tends to
        $1/2$, so the relaxation bound can be off by a factor of two and rounding is not
        a general method.
\end{enumerate}

\section*{Chapter 25: Neural Networks and Backpropagation}
\begin{enumerate}[label=\textbf{25.\arabic*}]
  \item The bound is $\lfloor1/\gamma^{2}\rfloor$: for $\gamma=0{,}2$ we get $25$,
        for $\gamma=0{,}05$ we get $400$. That bound only counts mistakes,
        and each mistake consumes margin, which is why $n$ does not appear. For
        small $\gamma$ the bound blows up like $1/\gamma^{2}$ and is therefore useless.

  \item Write $\mathbf{W}=\mathbf{W}_{2}\mathbf{W}_{1}$ and
        $\mathbf{b}=\mathbf{W}_{2}\mathbf{b}_{1}+\mathbf{b}_{2}$; then the output is
        $\mathbf{W}\mathbf{x}+\mathbf{b}$, that is, one linear threshold. Adding the next
        identity layer only multiplies matrices, so the capacity does not
        increase at all.

  \item With $\mathbf{x}=(0{,}5;-1;\,0{,}25;2)$ and $y=0$: $\mathbf{z}_{1}=(2{,}9125;\,
        0{,}375)$, $a=0{,}8961350845$, $\bar{z}_{2}=0{,}8961350845$, hence
        $\bar{\mathbf{z}}_{1}=(0{,}7169081;\,-0{,}5376810)$ and
        \[
          \nabla_{\mathbf{W}_{1}}L=
          \begin{pmatrix}
            0{,}3584540 & -0{,}7169081 & 0{,}1792270 & 1{,}4338161\\
            -0{,}2688405 & 0{,}5376810 & -0{,}1344203 & -1{,}0753621
          \end{pmatrix},
        \]
        that is, the second row equals $-0{,}75$ times the first row. This item is
        ambiguous: both gates are open, so there is no row with zero gradient.

  \item Not enough evidence. Finite differences only check thirteen scalars at one
        point, and at $h=10^{-6}$ an error of $10^{-10}$ is exactly what is expected; a
        shared bug in both computations would not show up. What is needed: several points,
        both gate states, the cases $y=0$ and $y=1$, gradients with respect to the input, and
        an automatic differentiation comparison.

  \item False. The theorem is about density of continuous functions on a
        compact domain: there exist parameters that approximate, not that those parameters
        are reachable by optimization. Two things are not stated: how many units are
        needed, and the error on test data (generalization), not only on the target
        function.

  \item At $t=1$, $\hat m=g$ and $\hat v=g^{2}$, so the step is
        $\eta g/(|g|+\epsilon)=-0{,}01$ and $b_{2}$ becomes $0{,}06$. With a gradient
        $100$ times larger the step stays $-0{,}01$, because Adam only
        uses the ratio to the gradient scale; $b_{2}$ stays $0{,}06$.

  \item 
The output $w_{2}\max(\mathbf{W}_{1}\mathbf{x}+\mathbf{b}_{1},0)$ is a sum
        of two linear functions, each active on one half-plane, hence piecewise
        linear; two lines $z_{1}[i]=0$ cut the plane into at most four cells, each
        with its own linear formula. For $k$ units in general position, $k$ lines
        cut the plane into $1+k+k(k-1)/2$ regions, that is $2,4,7$ for $k=1,2,3$
        (not $8$: the third line adds only three regions, not four).

  \item $\nabla_{\mathbf{x}}L=(-0{,}0868370;\,0{,}1328096;\,-0{,}0970532;\,
        -0{,}1430257)$, so $\boldsymbol{\delta}=0{,}05\,\mathrm{sign}(\nabla)
        =(-0{,}05;0{,}05;-0{,}05;-0{,}05)$. The loss at $\mathbf{x}+\boldsymbol{\delta}$
        is $0{,}3186800$, whereas the linear prediction
        $0{,}2949123+0{,}0229863=0{,}3178985$; a difference of $7{,}81\times10^{-4}$ from
        curvature.

  \item 
Prediction: the second lab block uses ReLU, and ReLU saturates only on the
        negative side. Saturation table ($8$ units, random He $\mathbf{W}_{1}$,
        $\mathbf{b}_{2}=0$):
        \begin{center}
        \begin{tabular}{lrrrr}
        & range of $Z_{1}$ (initial) & dead units & final $L$ & accuracy \\ \hline
        $\mathbf{b}_{1}=0$   & $[-2.43;\,2.40]$   & $0/8$ & $0.0163$ & $1.0000$ \\
        $\mathbf{b}_{1}=+10$ & $[7.57;\,12.41]$   & $0/8$ & $0.3225$ & $0.8375$ \\
        $\mathbf{b}_{1}=-10$ & $[-12.43;\,-7.60]$ & $8/8$ & $0.6909$ & $0.5333$
        \end{tabular}
        \end{center}
        So $\mathbf{b}_{1}=+10$ indeed does not kill the hidden units ($Z_{1}$ far
        above zero passes through unchanged), but it is not free: $H\cdot\mathbf{w}_{2}$
        becomes very large, the output sigmoid sticks at $0$ or $1$, and $L$ starts
        at $4.04$; after $3000$ steps it stops at $L=0.3225$ with accuracy $0.8375$,
        worse than $\mathbf{b}_{1}=0$. For $\mathbf{b}_{1}=-10$ every unit is dead
        from the start ($Z_{1}$ below zero, the derivative exactly $0$, not merely
        small), $\partial L/\partial b_{1}=\mathbf{0}$ exactly, and training stops at
        $L=0.6909$ -- the label entropy -- with majority-class accuracy $0.5333$.
  \end{enumerate}

\section*{Chapter 26: Generalization, Regularization, and Training}
\begin{enumerate}[label=\textbf{26.\arabic*}]
  \item Not enough evidence, and what is wrong is equating training error with
        generalization error. Zero training error only means the training data has been
        memorized. The additional numbers needed: validation (or test) error on data not
        used for training, together with its interval width.

  \item With $\hat f=\hat f_{D}$ and $\bar f=\mathbb{E}_{D}\hat f$:
        $\mathbb{E}(y-\hat f)^{2}=\sigma^{2}+(\bar f-f)^{2}+\operatorname{Var}(\hat f)$.
        The cross term vanishes because $\mathbb{E}[\varepsilon]=0$ and $\varepsilon$
        is independent of $D$, so $\mathbb{E}[\varepsilon(\hat f-f)]=0$; independence
        is used precisely at this step, on $\varepsilon$, not on $\hat f$. The second
        cross term vanishes for a different reason: $\mathbb{E}[\hat f-\bar f]=0$ by
        the definition of $\bar f$, so $2(f-\bar f)\,\mathbb{E}[\hat f-\bar f]=0$
        without needing any independence.

  \item From $\frac{d}{dw}\bigl[\frac12\sum(y_{i}-wx_{i})^{2}\bigr]=0$ we get
        $w(\lambda)=\sum x_{i}y_{i}/(\sum x_{i}^{2}+\lambda)=64/(30+\lambda)$.
        Its derivative $-64/(30+\lambda)^{2}<0$ but never zero, so $w$ shrinks
        without touching zero. In L1, the subgradient at $w=0$ contains zero when
        $\lambda\ge|\sum x_{i}y_{i}|=64$; so $w=0$ exactly at $\lambda=64$.

  \item $\operatorname{MSE}(c)=(c-1)^{2}\mu^{2}+c^{2}\sigma^{2}/n$: for
        $c=0{,}95$ we get $0{,}01+0{,}0361=0{,}0461$; for $c=1{,}05$ we get
        $0{,}01+0{,}0441=0{,}0541$; for $c=1$ we get $0{,}04$. An equally large deviation
        is not an equally large loss, because the variance term grows like $c^{2}$, so
        exceeding the budget costs more than shrinking.

  \item Fix: that cross-validation number is an estimate of the error of the selection
        procedure, not of the selected model's error in the field, and it is
        optimistically biased because $\lambda$ was chosen on the same folds. Two
        things are needed: nested cross-validation (or a test set never seen during
        selection), and repetition with several partitions along with intervals.

  \item For $a<0$: with $u<0$ the derivative is $u-a-t$, zero at $u=a+t$, which is valid
        only when $a<-t$. For $u=0$ the subgradient is $[-a-t,\,-a+t]$, which contains
        zero when $|a|\le t$, so $u=0$ is optimal. Symmetric because the objective
        $\frac12(u-a)^{2}+t|u|$ is unchanged when $u$ and $a$ are both sign-flipped.

  \item The half-life is $\ln2/(\eta\lambda)=0{,}693147/0{,}002=346{,}57$ steps;
        a quarter of the initial value means two half-lives, $693{,}15$ steps, and
        $(1-\eta\lambda)^{693}=0{,}24973$. A numpy simulation with zero gradient
        gives the same numbers, because $w_{k}=w_{0}(1-\eta\lambda)^{k}$.

  \item The gap is $0{,}2406-0{,}1500=0{,}0906$; at $\lambda=0$ the gap is
        $0{,}3262-0=0{,}3262$, at $\lambda=10$ only $0{,}3395-0{,}3333=0{,}0062$.
        A small gap is not enough: at $\lambda=10$ the gap is smallest precisely because
        both are equally bad, close to the majority-class error $0{,}3419$.

  \item With $m$ hypotheses, $P(\exists i:|\hat R_{i}-R_{i}|>\epsilon)\le
        2m e^{-2n\epsilon^{2}}$, so $\epsilon=\sqrt{(\ln m+\ln(2/\delta))/(2n)}$.
        For $m=2^{10}$, $n=200$, $\delta=0{,}01$: $\epsilon=0{,}1749$. Halving it
        demands four times the $n$, namely $800$.
\end{enumerate}

\section*{Chapter 27: Honest Model Evaluation}
\begin{enumerate}[label=\textbf{27.\arabic*}]
  \item Not enough evidence. A difference of $0{,}4$ points on a single test set is not
        necessarily beyond noise; what is wrong is treating two paired proportions
        as independent. Two additional numbers: the size of the test set, and the count of
        paired disagreements (for McNemar's test) or the interval width.

  \item For hard labels $s,y\in\{0,1\}$ we have $(s-y)^{2}=1$ exactly at
        mistakes and $0$ at correct predictions, so the average is
        $(\mathrm{FP}+\mathrm{FN})/n$. For scores that are probabilities, $(s-y)^{2}$
        takes intermediate values, so the average is no longer a count of mistakes
        but a calibration penalty.

  \item $\mathbb{E}[(s-y)^{2}]=s^{2}-2s\pi+\pi$ with $\pi=\mathbb{E}[y]$.
        The derivative is $2s-2\pi$, zero at $s=\pi$; the minimum value is
        $\pi^{2}-2\pi^{2}+\pi=\pi(1-\pi)$, and the second derivative $2>0$ confirms
        that it is a minimum.

  \item Precision $180/225=0{,}8$; recall $180/200=0{,}9$; $F_{1}=0{,}847059$;
        accuracy $1935/2000=0{,}9675$; specificity $1755/1800=0{,}975$; balanced accuracy
        $(0{,}9+0{,}975)/2=0{,}9375$. \texttt{sklearn.metrics} gives the
        same numbers to six decimals.

  \item Wilson $95/100$: $[0{,}8882;\,0{,}9785]$, width $0{,}0902$; for
        $950/1000$: $[0{,}9347;\,0{,}9619]$, width $0{,}0272$, that is $3{,}318$
        times narrower even though $n$ is ten times larger. Width is not a measure of
        model quality: it depends on $n$ and on the data, and the two experiments are
        not paired.

  \item 
Out of $25$ pairs, $19$ are correct: the positive $0{,}95$ beats five
        negatives, $0{,}88$ also five (it sits above $0{,}75$), $0{,}70$ four,
        $0{,}60$ three, $0{,}40$ two. Hence $\mathrm{AUC}=19/25=0{,}76$;
        \texttt{roc\_auc\_score} reports $0{,}76$.

  \item Precision $=0{,}8\pi/(0{,}8\pi+0{,}01(1-\pi))$: at $\pi=0{,}002$ we get
        $0{,}138169$, at $\pi=0{,}2$ we get $0{,}952381$. The requirement precision
        $\ge0{,}5$ gives $0{,}4\pi\ge0{,}005-0{,}005\pi$, that is
        $\pi\ge0{,}005/0{,}405=0{,}012346$.

  \item Split by host, not by record: $60/20/20$ becomes $180$ training hosts,
        $60$ validation, $60$ test, each about $20$ records. Preprocessing is
        computed inside the training fold only. The choice among $40$ settings goes
        through nested cross-validation in the training part. The test set is touched
        once, at the end.

  \item False. AUC measures ranking, not usefulness at a particular threshold. With
        sensitivity $0{,}8$ and specificity $0{,}99$, at prevalence $0{,}002$ the
        precision is only $0{,}138169$: of ten alarms, $8{,}62$ are wrong, so about
        nine. The chosen
        threshold still has to be checked for precision and recall.
\end{enumerate}

\section*{Chapter 28: Formal Methods with Lean 4}
\begin{enumerate}[label=\textbf{28.\arabic*}]
  \item False. A file that is accepted means every theorem has a proof that is
        checked, but \texttt{sorry} passes with a warning, and
        \texttt{axiom} as well as \texttt{native\_decide} extend the trusted base.
        The checking command: \texttt{\#print axioms} on every theorem that is
        claimed.

  \item 
Helper lemma: for every $k$, $k(2k)+2(2k+1)=(k+1)\cdot2(k+1)$. That goal
        contains $k\cdot k$, so it is not a linear goal: \texttt{omega} will not
        close it on its own and a rewrite is needed first. Then \texttt{induction n
        with}: base case \texttt{rfl}; the \texttt{succ k ih} step uses \texttt{rw
        [oddSum, Nat.mul\_add, ih]}, which leaves exactly that helper lemma, then
        \texttt{exact bantu k}. The helper lemma itself is closed by
        \texttt{simp only [Nat.mul\_add, Nat.add\_mul, Nat.mul\_one, Nat.one\_mul]}
        followed by \texttt{omega}: \texttt{omega} cannot close
        $k\cdot(2k)+2(2k+1)=(k+1)\cdot2(k+1)$ on its own, because that goal contains
        $k\cdot k$ and is not linear (tried: \texttt{omega could not prove the goal}).

  \item Informal: for $n=0$ both sides are $0$; if $\sum_{i<n}2i=n(n-1)$, then
        adding $2n$ gives $n(n-1)+2n=(n+1)n$. The Lean version defines
        \texttt{evenSum : Nat -> Nat} recursively (\texttt{evenSum 0 = 0},
        \texttt{evenSum (n+1) = evenSum n + 2*n}) and states the form without
        truncated subtraction, $E(n)+n=n\cdot n$. The steps: \texttt{induction},
        base \texttt{rfl}, then \texttt{rw [evenSum]}. Note that
        \texttt{rw [evenSum, ih]} fails outright (\texttt{did not find instance of
        the pattern evenSum k + k}), because the goal afterwards reads
        $E(k)+2k+(k+1)=(k+1)(k+1)$. Insert the regrouping
        \texttt{have h : evenSum k + 2 * k + (k + 1) = (evenSum k + k) + (2 * k + 1) := by omega}
        first, continue with \texttt{rw [h, ih]}, and close with
        \texttt{simp only [Nat.mul\_add, Nat.add\_mul, Nat.mul\_one, Nat.one\_mul]}
        then \texttt{omega}.

  \item The sum of the first $37$ odd numbers $=37^{2}=1369$; the sum of the first $250$
        $=250^{2}=62500$. Script: \texttt{sum(2*i+1 for i in range(n))} gives
        both numbers, and both are perfect squares.

  \item        A direct check for all $n\le1000$ demands $1001$ cases, namely $n=0$
        through $n=1000$; the bound $n\ge1$ would give $1000$, but the exercise
        does not say that, and the \texttt{zero} base case above really does check
        $n=0$. An induction proof generates only two goals, \texttt{zero} and
        \texttt{succ}, then concludes for all $n$; so $1001$ against $2$.

  \item \texttt{Nat.add} is defined by recursion on the second argument, so
        $n+0$ reduces straight to $n$ and \texttt{rfl} closes it. On $0+n$
        the recursion gets stuck on the unknown $n$, so there is nothing
        to reduce and \texttt{induction n} is needed.

  \item 
Weak: for fixed $x$, $\exists\delta,\ |\delta|\le1\wedge f(x+\delta)=f(x)$;
        the witness $\delta=0$ makes both sides equal and \texttt{rfl} closes it.
        Strong: $\forall x\,\forall\delta,\ |\delta|\le1\to f(x+\delta)=f(x)$, and
        that is false for the identity $f$ at $\delta=1$.

  \item 
\texttt{decide}: no axioms, pure kernel computation. \texttt{omega}:
        \texttt{propext} and \texttt{Quot.sound}, because it proves goals by
        contradiction inside the integers; \texttt{Classical.choice} does not
        appear, and this chapter prints the same list. \texttt{sorry}:
        \texttt{sorryAx}, the axiom added specifically to close a skipped goal.

  \item First mistake: \texttt{delta} is already free in the context and then bound again
        by \texttt{exists}, so the reader of the statement determines its value,
        not the attacker; robustness is not implied. Second: there is no quantifier over $x$,
        so $x$ is free too. Correct version:
        $\forall x,\ \exists\delta,\ \lVert\delta\rVert\le1\wedge f(x+\delta)=f(x)$.

  \item \texttt{intro h} gives $h:P\wedge\neg P$. Its components \texttt{h.1} have type
        $P$ and \texttt{h.2} have type $\neg P$; then \texttt{exact h.2 h.1} closes
        the goal, because $\neg P$ is $P\to$ false.
\end{enumerate}

\section*{Chapter 29: Adversarial Machine Learning: Attacks and Defenses}
\begin{enumerate}[label=\textbf{29.\arabic*}]
  \item $\mathbf{s}_{A}=1$, $\mathbf{s}_{B}=3$, so the margin is already $-2$ (the example
        leans toward $B$); $\mathbf{d}=(1;-2;1)$ with $\lVert\mathbf{d}\rVert_{1}=4$
        and $\lVert\mathbf{d}\rVert_{2}=\sqrt6\approx2{,}449$. Hence
        $\varepsilon^{*}_{\infty}=2/4=0{,}5$ and
        $\varepsilon^{*}_{2}=2/2{,}449=0{,}8165$; the ratio is $\sqrt{8/3}=1{,}633$.

  \item The derivative of the objective with respect to $\boldsymbol{\delta}$ is $\mathbf{w}$, so
        the maximum is $\varepsilon\lVert\mathbf{w}\rVert_{2}$, attained uniquely at
        $\boldsymbol{\delta}=\varepsilon\mathbf{w}/\lVert\mathbf{w}\rVert_{2}$
        (Cauchy--Schwarz, with equality when parallel). In
        $\ell_\infty$ the maximum is $\varepsilon\lVert\mathbf{w}\rVert_{1}$ at
        $\varepsilon\,\mathrm{sign}(\mathbf{w})$, a vertex; the form differs
        because the dual norm differs.

  \item The error $\mathbf{d}^{\top}\boldsymbol{\delta}$ is linear on a compact convex
        set, so the maximum is at an extreme point. The extreme points of the
        $\ell_\infty$ ball are $\pm\varepsilon$ in each coordinate, and FGSM picks
        exactly that point; hence the PGD path, which always stays inside the ball,
        cannot exceed it.

  \item False. In this chapter a random signed perturbation at $\varepsilon=0{,}15$ only
        drops accuracy to $95{,}48$ percent, whereas PGD 40 steps
        brings it down to $10{,}72$ percent. Robustness to random noise
        says nothing about an attack that chooses the direction.

  \item 
Untargeted success (PGD with 40 steps, clean accuracy as the reference) for
        $A$: $14{,}02$, $53{,}16$, $89{,}06$ percent at $\varepsilon=
        0{,}05;0{,}10;0{,}15$. For $D$: $5{,}38$, $15{,}28$, $35{,}06$ percent. Over
        that range the gap actually widens ($8{,}64$, $37{,}88$, $54{,}00$
        percentage points), because $A$ is already close to the $100$ percent
        ceiling while $D$ is still far; the narrowing only happens outside that
        range, at $\varepsilon\ge0{,}20$, once both models run out of margin.

  \item One coordinate suffices: projection onto an interval is clipping, and
        $[x_{j}-\varepsilon,x_{j}+\varepsilon]\cap[0,1]$ equals
        $[\max(0,x_{j}-\varepsilon),\ \min(1,x_{j}+\varepsilon)]$. Clipping to the
        $\varepsilon$ box and then to $[0,1]$, or the other way round, yields the
        same interval because interval intersection is commutative; the Cartesian product
        then gives the same point.

  \item $D$ loses $97{,}99-96{,}48=1{,}51$ points of clean accuracy to
        gain $62{,}65-10{,}72=51{,}93$ points of robustness, so the price is
        $1{,}51/51{,}93=0{,}0291$ clean points per point of robustness. That ratio is not
        complete: it depends on one $\varepsilon$, on one attacker that is
        not an upper bound, and it ignores the clean accuracy of the comparison model.

  \item Protocol: attack with PGD with more steps and several random starts,
        with EOT over the defense's randomness, then repeat with attacks without
        gradients (imitation/transfer, finite-difference estimation). Report the curve of accuracy
        against $\varepsilon$, not a single point. Easiest symptom: success
        jumps as soon as the gradient is replaced by BPDA or EOT, while naive attacks
        report zero.

  \item White box $89{,}06$ percent, transfer $1-0{,}2379/0{,}9849=75{,}85$
        percent, a difference of $13{,}21$ points. That is not evidence that $B$ is safer:
        an attack built specifically for $B$ will be stronger, and in this chapter the
        white-box versus transfer difference does shrink when the attacker is adapted.
        This item is ambiguous: the lab output does not contain a PGD row for model $B$.

  \item Untargeted only needs a label change; targeted must reach a particular
        class, so the set of targets that must be met is far smaller. At
        $\varepsilon=0{,}15$ the target is reached $33{,}00$ percent of the time, whereas
        untargeted is $89{,}06$ percent; with $\varepsilon=0{,}30$ the target is reached
        $99{,}50$ percent. So the difficulty is a matter of budget, not of attack type.
\end{enumerate}

\section*{Chapter 30: The Limits of Attack: Obfuscated Gradients and NIDS Evasion}
\begin{enumerate}[label=\textbf{30.\arabic*}]
  \item Three checks: (a) repeat with adaptive attacks (BPDA, EOT), (b)
        report the curve of accuracy against $\varepsilon$, not a single point, (c) check
        gradient quality (cosine against finite differences or against the weight direction). If
        what is happening is gradient masking: success collapses only for attacks
        based on gradients, while random search and transfer still succeed, and
        the curve is flat.

  \item For $f(x)=w\cdot x+b$ the gradient is constant and the Hessian zero, so the
        Lipschitz constant is $L=0$ and the bound $L\delta$ gives zero error: BPDA is exact. With two
        ReLU layers, $f$ is piecewise linear so the gradient jumps at the
        kinks and the Lipschitz-gradient hypothesis fails; the bound is then
        meaningless.

  \item 
As the program shows, the label flips up to $q=0{,}6$ (score $-1{,}2$) and
        fails at $q\ge0{,}70$, when every coordinate rounds to zero so the score is
        $0$, a tie rather than a flip. Hence
        $q^{*}=2\max_{j}|x_{j}-\varepsilon|=2(\varepsilon+|\min_{j}x_{j}|)=0{,}70$.
        Analytic safety bound: the rounding error is at most $4q$, so the attack is
        only guaranteed to succeed when $q<(8\varepsilon-m)/4=0{,}345$; at $q=0{,}345$
        the label still flips, so the direction of the inequality matters.

  \item Probability of rejection $=P(\text{lebih dari }10\text{ dari }20\text{ hilang})
        =0{,}411901$, so what is answered is $588$ of $1000$ queries: $588/20=29$
        PGD steps with $m=20$. With $m=1$ the rejection probability is $0{,}5$, so $500$
        steps. What is paid: the gradient becomes noisier and its direction further from the true one.

  \item What can be predicted without running: clean accuracy is about the same,
        naive success stays near zero, and BPDA success is near the baseline for $q=0$
        and $q=0{,}05$ but falls for large $q$. The cosine no longer approaches one
        because the ReLU gate adds kinks on top of rounding, so the BPDA approximation
        only approaches the direction. The numbers for the third value of $q$ must be taken from execution,
        not invented.

  \item Binary flags and protocol numbers are assumed not to be changeable by the attacker:
        $7+3=10$ locked features, the remaining $40$ features can really be shifted (out of
        $50$). A $32\times32$ image with three channels has $3072$ values, so
        $3072/40=76{,}8$ times more.

  \item Daily retraining changes the parameters that a gradient-based attack uses,
        so old directions go stale quickly; a blind attack does not use those parameters, so
        it keeps working. According to the EOT lemma what must be checked is the expectation over
        transformations, and an adaptive attack must be defined against that defense
        itself, including its retraining schedule.

  \item Metric: success of the adaptive attack against $\varepsilon$, with
        at least $5$ seeds and an interval. The curve shape that accepts (a): the curve
        is shifted to the right or downward systematically at all $\varepsilon$,
        not only at one point. If after BPDA and EOT the success returns
        to the baseline, then only the measurement was broken, so (b).
\end{enumerate}

\section*{Chapter 31: Cryptography: Modular Integers to the Quantum Threat}
\begin{enumerate}[label=\textbf{31.\arabic*}]
  \item False. $2^{100}\equiv1\pmod{101}$ only says the order of $2$ divides $100$,
        that is, $101$ passes the Fermat test to base $2$; pseudoprimes (and
        Carmichael numbers) also pass. What that congruence proves is only $101\nmid2$
        and $2^{100}\equiv1$, not the absence of divisors.

  \item If $\gcd(a,n)=\gcd(b,n)=1$, Bezout's theorem gives $ax+ny=1$ and
        $bz+nw=1$; the product of the two gives $(ab)(xz)+n(\cdot)=1$, so $ab$ is a unit
        and $U(n)$ is closed. If $a$ is not a unit, $d=\gcd(a,n)>1$ and
        $b=n/d\ne0$ satisfies $ab=n(a/d)\equiv0$.

  \item $a^{2}-1=(a-1)(a+1)\equiv0\pmod p$; since $p$ is prime, Euclid's lemma gives
        $p\mid a-1$ or $p\mid a+1$, so $a\equiv\pm1$. For modulus $8$ that lemma
        fails because $8$ is not prime: the square roots of one are $1,3,5,7$, since
        $3^{2}=5^{2}=7^{2}=49\equiv1$.

  \item The classes $0\le a<pq$ that share a factor with $pq$ are $q$ multiples of $p$
        plus $p$ multiples of $q$, and $0$ is counted twice; so
        $\varphi(pq)=pq-p-q+1=(p-1)(q-1)$. For $p=1019$, $q=1031$:
        $\varphi=1018\cdot1030=1\,048\,540$.

  \item Extending Euclid backwards gives $17^{-1}\equiv2753\pmod{3120}$ (since
        $17\cdot2753=46\,801=15\cdot3120+1$), and
        $65\,537^{-1}\equiv60\,493\pmod{1\,048\,540}$. The division takes $4$ steps against $5$,
        so adding one digit to the modulus adds only a few
        steps.

  \item Fermat's little theorem: $a^{100}\equiv1\pmod{101}$ for $101\nmid a$. Hence
        $5^{100}\equiv1$, $3^{200}=(3^{100})^{2}\equiv1$, and
        $2^{1000}=(2^{100})^{10}\equiv1$. All three are checked by
        \texttt{pow(a,k,101)} and match.

  \item $A=5^{7}\equiv17$, $B=5^{15}\equiv19$, shared key $19^{7}\equiv15$
        (the same as $17^{15}$). If an eavesdropper replaces $A$ with $11$, Bob
        computes $11^{15}\equiv10$ while Alice still holds $15$; two different
        keys appear because the paired exponents are no longer a legitimate
        Diffie--Hellman pair.

  \item With $e=3$ and $m<n^{1/3}$ we have $m^{3}<n$, so $c=m^{3}$ with no
        modular reduction and the integer cube root returns $m$.
        For $n=1\,050\,589$ the bound is $n^{1/3}\approx101{,}66$, so the longest
        message is $m=101$. Random padding makes $m^{3}>n$ so that integer root
        does not exist.

  \item Two things must be true: (i) the parameters and implementation really
        meet the claimed security guarantee, that is, the reduction to a lattice
        problem that is still hard; (ii) the hardness assumption underneath still
        holds and the proof model covers the attacks that are actually used.
        Widespread adoption is not a condition for the claim to be true: a scheme
        used everywhere is still not proven secure if its assumption fails.
        The second sentence remains wrong: RSA is still embedded in legacy systems,
        and migration is done hybrid, so it takes years. (Post-quantum signatures
        already exist: ML-DSA, FIPS 204.)

  \item About $1{,}1774\sqrt{2^{48}}=1{,}1774\cdot2^{24}$, that is $19\,753\,662$
        samples for a one-half collision probability; for $32$ bits
        $1{,}1774\cdot2^{16}=77\,163$. The numbers $2^{24}$ and $2^{16}$ are rough
        square-root guesses, not the exact values.
\end{enumerate}

\section*{Chapter 32: Differential Privacy and the Mathematical Foundations of QML}
\begin{enumerate}[label=\textbf{32.\arabic*}]
  \item False as a sentence. $\varepsilon=0{,}5$ bounds the probability ratio per
        record, it does not remove leakage; it is also a worst-case guarantee. Two things
        that must be true: the mechanism really satisfies $(\varepsilon,\delta)$ for
        the adjacency relation and sensitivity used, and all outputs
        (queries, budget remainder, post-processing) are counted within that budget.

  \item The $d$-dimensional Laplace density is a product, so the ratio is
        $\prod_{j}\exp\bigl((|y_{j}-a'_{j}|-|y_{j}-a_{j}|)/b\bigr)$, and by the
        triangle inequality each factor is at most
        $\exp(|a_{j}-a'_{j}|/b)$. The product is
        $\exp(\sum_{j}|a_{j}-a'_{j}|/b)\le\exp(\Delta f/b)=e^{\varepsilon}$.
        The one-dimensional argument is replaced by the product of per-coordinate ratios, and
        $\ell_1$ takes the place of the single absolute value.

  \item Take $D\sim D'$ differing in one record. That record exists in only
        one of $D_{1},D_{2}$, so only one mechanism has its input changed;
        its ratio is bounded by $e^{\varepsilon_{i}}$, while the others have ratio $1$.
        Hence $\varepsilon=\max_{i}\varepsilon_{i}$ and
        $\delta=\max_{i}\delta_{i}$, not the sum.

  \item The probability of answering ``yes'' is $p+(1-p)/2$ if the truth is ``yes'' and $(1-p)/2$
        if the truth is ``no''; the ratio is $(1+p)/(1-p)$. Hence
        $p=(e^{\varepsilon}-1)/(e^{\varepsilon}+1)$, and for $\varepsilon=2$
        we get $p=0{,}761594$, that is $\tanh1$.

  \item Count query: $b=2$, so $P(|Y|>10)=e^{-5}=0{,}006738$. Mean query
        with $R=1$: $b=10^{-4}/0{,}5=2\times10^{-4}$; an error of $10$ people means
        $10^{-3}$ on the mean scale, which again gives $e^{-5}$, whereas an
        absolute error of $10$ on that scale gives $e^{-50000}\approx0$. This item is
        ambiguous about the units of the error.

  \item $H|0\rangle=(|0\rangle+|1\rangle)/\sqrt2$; $S$ gives
        $(|0\rangle+i|1\rangle)/\sqrt2$; then $H$ gives
        $\frac{1+i}{2}|0\rangle+\frac{1-i}{2}|1\rangle$. The probabilities are $1/2$ and $1/2$,
        the same as $H|0\rangle$, because measurement in the computational basis only sees the modulus
        of the amplitude. What distinguishes them is the relative phase, and the gate is $S^{\dagger}$
        (it multiplies $|1\rangle$ by $-i$). $Z$ after $H$ is no substitute: $HZH=X$,
        which swaps $|0\rangle$ with $|1\rangle$ instead of rotating the phase.

  \item $|\Phi^{+}\rangle=(|00\rangle+|11\rangle)/\sqrt2$. If the second qubit is measured
        first and the result is $0$, the first qubit collapses to $|0\rangle$, so the
        probability of the outcome $01$ is $0$; if the result is $1$, the first qubit collapses
        to $|1\rangle$, so the outcome $11$ is certain and the probability of $01$ is also $0$.
        Both show perfect correlation: the outcomes of the two qubits are always the same.

  \item $U^{\dagger}U=\operatorname{diag}(1,e^{-i\varphi})
        \operatorname{diag}(1,e^{i\varphi})=I$, so it is unitary for every
        $\varphi$. In the computational basis $|1\rangle$ only picks up a phase, its
        amplitude modulus stays the same, so probabilities do not change; but that relative phase
        changes the interference at the next $H$ gate, so the result can
        differ.

  \item 
$n=40$: $2^{40}=1{,}0995\times10^{12}$ amplitudes, times $16$ bytes, gives
        $16$ TiB (not GiB: $2^{40}$ bytes is already one TiB). $n=60$:
        $2^{60}=1{,}1529\times10^{18}$ amplitudes, that is $16$ EiB or
        $1{,}8447\times10^{19}$ bytes, far beyond any storage that exists.

  \item Three questions: whether all the amplitudes are really used, whether
        the output can be extracted with a number of measurements that is not exponential,
        and whether there is a classical comparison on the same task and data. What makes
        me accept: a proven separation or measurable scaling on that task,
        a polynomial number of measurements, and an advantage over a seriously tuned
        classical baseline.
\end{enumerate}
